\Askhsh{1778}{4}
\begin{ylh}{1}
    \item 3.6. Κριτήρια ισότητας ορθογώνιων τριγώνων
    \item 5.9. Μια ιδιότητα του ορθογώνιου τριγώνου
    \item 5.10. Τραπέζιο.
\end{ylh}
Δίνεται ορθή γωνία $x\hat{\mathrm{O}}y$ και τα σημεία $\mathrm{A}$ και $\mathrm{B}$ των ημιευθειών $\mathrm{O}y$ και $\mathrm{O}x$ αντίστοιχα με $\mathrm{O}\mathrm{A} = \mathrm{O}\mathrm{B}$. Μία ευθεία $(\varepsilon)$ η οποία δεν είναι παράλληλη στην $\mathrm{A}\mathrm{B}$ διέρχεται από το $\mathrm{O}$ ώστε τα σημεία $\mathrm{A}$ και $\mathrm{B}$ να είναι στο ίδιο ημιεπίπεδο. Η κάθετη από το $\mathrm{A}$ στην $(\varepsilon)$ την τέμνει στο $\Delta$ και η κάθετη από το $\mathrm{B}$ στην $(\varepsilon)$ την τέμνει στο $\mathrm{E}$.
Να αποδείξετε ότι:
\begin{alist}
    \item Τα τρίγωνα $\mathrm{O}\mathrm{A}\Delta$ και $\mathrm{O}\mathrm{E}\mathrm{B}$ είναι ίσα. \monades{7}
    \item $\mathrm{A}\Delta + \mathrm{B}\mathrm{E} = \Delta\mathrm{E}$. \monades{7}
    \item $\mathrm{M}\mathrm{N} = \frac{\Delta\mathrm{E}}{2}$, όπου $\mathrm{M}$ και $\mathrm{N}$ τα μέσα των $\mathrm{A}\mathrm{B}$ και $\Delta\mathrm{E}$ αντίστοιχα. \monades{7}
    \item Το τρίγωνο $\Delta\mathrm{M}\mathrm{E}$ είναι ορθογώνιο και ισοσκελές. \monades{4}
\end{alist}
